% ============================================================
\part{Docker}
% ============================================================

\chapter{Wir wollen ein Programm auf jedem Computer gleich ausführen}

\kapitelziel{
Dieses Kapitel wird nach dem Werkstatt-Schema ausgearbeitet.
}

\kapitelstruktur

\begin{aufgabeBox}
Aufgabe folgt.
\end{aufgabeBox}

% TODO: Situation, Problem, Wissen, Lösung, Merksatz,
%       Fehler, Übungen, Quiz, Haube, Robotik-Transfer

\chapter{Wir wollen einen Linux-Container betreten}

\kapitelziel{
Dieses Kapitel wird nach dem Werkstatt-Schema ausgearbeitet.
}

\kapitelstruktur

\begin{aufgabeBox}
Aufgabe folgt.
\end{aufgabeBox}

% TODO: Situation, Problem, Wissen, Lösung, Merksatz,
%       Fehler, Übungen, Quiz, Haube, Robotik-Transfer

\chapter{Wir wollen Daten behalten, obwohl der Container gelöscht wird}

\kapitelziel{
Dieses Kapitel wird nach dem Werkstatt-Schema ausgearbeitet.
}

\kapitelstruktur

\begin{aufgabeBox}
Aufgabe folgt.
\end{aufgabeBox}

% TODO: Situation, Problem, Wissen, Lösung, Merksatz,
%       Fehler, Übungen, Quiz, Haube, Robotik-Transfer

\chapter{Wir wollen unser eigenes Docker-Image bauen}

\kapitelziel{
Dieses Kapitel wird nach dem Werkstatt-Schema ausgearbeitet.
}

\kapitelstruktur

\begin{aufgabeBox}
Aufgabe folgt.
\end{aufgabeBox}

% TODO: Situation, Problem, Wissen, Lösung, Merksatz,
%       Fehler, Übungen, Quiz, Haube, Robotik-Transfer

\chapter{Wir wollen mehrere Container gemeinsam starten}

\kapitelziel{
Dieses Kapitel wird nach dem Werkstatt-Schema ausgearbeitet.
}

\kapitelstruktur

\begin{aufgabeBox}
Aufgabe folgt.
\end{aufgabeBox}

% TODO: Situation, Problem, Wissen, Lösung, Merksatz,
%       Fehler, Übungen, Quiz, Haube, Robotik-Transfer

